% Assume-se que \textual já foi feito

\chapter{Exemplo}
\label{ch:exemplo}

\index{bobagem} Primeiro parágrafo da seção com uma frase sem sentido que só serve para ocasionar uma quebra e de demonstrar a configuração de indentação da primeira linha. Essa frase está aqui pois parágrafos de uma linha são feios.

Resultado do uso de siglas:
\begin{itemize}
\item Sigla que nunca expande: \API;
\item Sigla normal, expande no primeiro uso: \DHT, mas não no segundo: \DHT;
\item Siglas com plurais automaticos: \APIs e \DHTs;
\item Plural não-trivial: \SQs;
\item Forçando uma expansão (e no plural) \Glsfirstplural{DHT};
\item Usando uma sigla cujo comando é diferente da sigla: \WTC.
\end{itemize}

Resultado do glossário:
\begin{itemize}
\item Dois termos, \polling e \proxy;
\item Plural: \proxys.
\end{itemize}

Resultado de \mla|index|: primeiro um link normal \indexterm{tomate}, depois um capitalizado \indexTerm{tomate}.

Exemplo de autoref com \mla|amsthm|: \autoref{def:definition}, \autoref{def:defn}, \autoref{lemma:1}, \autoref{theorem:1}, \autoref{proof:1}.

\begin{definition}
  \label{def:definition}
  Exemplo de definição.
\end{definition}

\begin{defn}
  \label{def:defn}
  Exemplo de definição com ambiente defn.
\end{defn}

\begin{lemma}
  \label{lemma:1}
  Exemplo de Lema.
\end{lemma}

\begin{theorem}
  \label{theorem:1}
  Exemplo de teorema.
  \begin{proof}
    Exemplo de prova dentro do teorema \qed
  \end{proof}
\end{theorem}

\begin{theoremproof}
  \label{proof:1}
  Exemplo de prova \qed
\end{theoremproof}

(Sub)enumerações e citações (verificar se OK com o idioma):
\begin{enumerate}
\item \cite{turing1937}:
  \begin{enumerate}
  \item \citeonline{turing1937}:
    \begin{enumerate}
    \item \citeonline{dijkstra1968};
    \end{enumerate}
  \end{enumerate}
\item \cite{turing1937,dijkstra1968};
\item \citeonline{turing1937,dijkstra1968}.
\end{enumerate}


\begin{listing}[tb]
\caption{Meta informações do presente documento.}
\label{lst:meta}
\begin{minted}[highlightlines={1,4-5}]{latex}
\titulo{Template \LaTeX{} para testes e dissertações do LAPESD/UFSC}
\autor{André Roberto Alves de Oliveira}
\data{1 de agosto de 2019} % ou \today
\tese % ou \dissertacao
\titulode{Mestre em Ciência da Computação}
\orientador{Prof. Frank Augusto Siqueira, Dr.}
% \coorientador{Prof. Lars Thørväld, Dr.}

\membrobanca{Prof. Valerie Béranger, Dr.}{Universidade Federal de Santa Catarina}
\membrobanca{Prof. Mordecai Malignatus, Dr.}{Universidade Federal de Santa Catarina}
\membrobanca{Prof. Huifen Chan, Dr.}{Universidade Federal de Santa Catarina}
\coordenador{Prof. Charles Palmer, Dr.}
\end{minted}
\fonte{o autor.}
\end{listing}


Resultado de \mla|\autoref|s:
\begin{itemize}
\item \autoref{lst:meta};
\item \autoref{alg:algoritmo};
\item \autoref{fig:figura} tem subfiguras:
  \begin{itemize}
  \item \autoref{fig:svg}
  \item \autoref{fig:brasao}
  \end{itemize}
\item \autoref{tb:tabela};
\item \autoref{ch:exemplo};
\item \autoref{sec:frutas};
\item \autoref{sec:goiaba};
\item \autoref{sec:jabuticaba};
\item \autoref{sec:tomate}.
\end{itemize}


\begin{algorithm}
  \caption{Exemplo do ambiente \texttt{algorithimic}.}
  \label{alg:algoritmo}
  \begin{algorithmic}[1]
    \Procedure{Closure}{C, A}
      \State{$H \gets \emptyset$}\Comment{Direct cache}
      \For{$i \in [1, n]$}\Comment{Parallel, (dynamic,32) scheduling}
        \State{$H \gets H \cup \Call{DoImportantStuff}{i}$}
      \EndFor
    \EndProcedure
  \end{algorithmic}
  \fonte{o autor.}
\end{algorithm}

\begin{figure}[tb]
  \centering
  \caption{Exemplo de figura com duas subfiguras.}   
  \label{fig:figura}
  
  % Subfiguras são feitas usando as funcionalidades do memoir. Não
  % inclua outros pacotes, pois eles podem fazer o memoir dar ragequit
  % 
  % Há duas maneiras, a maneira limpinha (só no lapesd-thesis.cls) e a
  % maneira do memoir (aviso: \subtop não funciona direito).
  \subcaptionminipage[fig:svg]%
    {.49\linewidth}%
    {O Makefile compila SVGs em PDFs usando o inkscape}%
    % {\includegraphics[width=.2\linewidth]{alphachannel.pdf}}%
  \hfill% 
  % o comando acima expande para o equivalente disso:
  \begin{minipage}[t]{.49\linewidth}%
    \centering
    \subcaption{Brasão da UFSC.\label{fig:brasao}}
    \includegraphics[width=.2\linewidth]{\jobname-logo.pdf}
  \end{minipage}

  \fonte{o autor.}
\end{figure}

\begin{table}[tb]
  \centering
  \caption{Exemplo de tabela e símbolos}
  \label{tb:tabela}
  \begin{tabular}{lccp{5cm}}
    \toprule
    Esquerda & Coluna 1    & \rotatebox{90}{90 graus}  & Parágrafo com \mla|p{5cm}|   \\
    \midrule
    $r_1$    & \cmk        &  \xmk                     & \circledi    \\
    $r_2$    &     \multicolumn{2}{c}{merged cell}     & \circledii   \\
    $r_3$    & \circlediii & \circlediv                & \circledv    \\
    $r_4$    & \circledvi  & \circledvii               & \circledviii \\
    $r_5$    & \circledix  &  x                        & y           \\
    \bottomrule 
  \end{tabular}
  \fonte{o autor.}
\end{table}

\section{Introdução}
\subsection{Contextualização do Problema}
\subsection{Objetivos}  
\subsection{Proposta}
\subsection{Contribuições}
\subsection{Organização do Texto}

\section{Fundamentação Teórica}
\subsection{Gêmeo Digital}
\subsubsection{Conceitos}

The concept of Digital Twins (DT) originated from Michael Grieves' work on Product Lifecycle Management (PLM),
the original term was ``Conceptual Ideal for PLM", later he referred to as the ``Mirrored Spaces Model" and the conceptual model
was called ``Information Mirroring Model", the term currently used, Digital Twin, was coined at NASA \cite{Grieves2017}.\\
The first use of Digital Twin technology, considered by the academia, originated at NASA during the Apollo 13 incident.
The need to simulate the effects of the anomaly to understand the problem and find a way to bring the astronauts back safely,
through real-time testing to evaluate solutions, led the engineers in charge to modify the pilot training simulators to replicate
the current state of the lunar module \cite{Allen2021}.\\
Formally, a DT is defined as a virtual replica of a physical asset, process, or system, being a software composed of models 
and services, that intentionally, represent the original system throught its lifecycle \cite{Dalibor2022}.

\subsubsection{Classificação}

The concept of Digital Twins in the literature often suffers from a lack of distinction between its representations and different
levels of complexity. Although the concept of DTs was originally introduced by Grieves in 2002, the technical differentiation between
variations of virtual replicas was consolidated by \citeonline{Kritzinger2018}. As proposed by \citeonline{Kritzinger2018}, the differentiation
of the digital model occurs through the level of integration of the data flow between the physical and virtual worlds.
The representation of the models is categorized into three different levels: Digital Twin, Digital Model, and Digital Shadow.\\
The differentiation proposed by \citeonline{Kritzinger2018} establishes an information flow as the dividing line between the categories.
While manual flow characterizes the initial models, increasing automation and the way data is exchanged define the transition to the
terms ``shadow'' and ``digital''. This distinction is illustrated in the communication logic between physical and virtual assets,
as detailed below:
\begin{itemize}
    \item Digital Model
    \item Digital Shadow
    \item Digital Twin
\end{itemize}
%A diferença entre Digital Model, Digital Shadow e Digital Twin%
% Talvez seja interessante colocar a definição de 5 camadas do FULLER aqui tbm %

\subsection{Arquitetura de Referência e Ciclo de Vida}

%Falar desde o começo do ciclo sobre os sensores, depois dados e sobre os serviços e a interface%

\subsection{Tecnologias Habilitadoras}

%breve resumo sobre os enablers%

\subsection{Do Gêmeo Único ao Ecossistema}

%Citar os trabalhos que falam sobre a execução de multiplos DTs%

\subsection{Desafios de Execuções e Escalabilidade}

%falar dos problemas em outras pesquisas%
\section{Trabalhos Relacionados / Revisão da Literatura}

\section{Proposta}
\section{Metodologia}
\section{Cronograma}

\section{Frutas}
\label{sec:frutas}
\lipsum[4]

\subsection{Goiaba}
\label{sec:goiaba}
\lipsum[4]

\subsubsection{Jabuticaba}
\label{sec:jabuticaba}
\lipsum[4]

\subsubsubsection{Tomate}
\label{sec:tomate}

\begin{figure}[tb]
  \centering
  \caption{Segunda Figura.}
  \label{fig:segunda-fig}
  \includegraphics[width=.2\linewidth]{\jobname-logo.pdf}
  \fonte{o autor.}
\end{figure}

\xindex{tomate} \lipsum[4]


%%% Local Variables:
%%% mode: latex
%%% TeX-master: "main"
%%% End:
